
\chapter{Reproducibility record}
\label{app:reproducibility}

This appendix records the configuration values that the main chapters rely on but do not
state, so that a reader can interpret the reported numbers and repeat the experiments. It
holds the prompt texts, question sets, settings, artifact identifiers, a worked prompt example
and a glossary. The measurements themselves are in
Chapters~\ref{ch:memory_evaluation} to~\ref{ch:spe} and the statistical detail is in
Appendix~\ref{app:on_device_statistics}.

\section{Prompts}
\label{app:repro_prompts}

Each prompt below was held fixed across every condition in which it was used, so that a
difference between conditions cannot be a difference in instructions. The texts are reproduced
from the harness source. Long lines are wrapped for display, and placeholders in braces are
filled at run time.

\subsection*{Reader answer prompt (\texttt{qa\_short\_answer\_v3})}

Used for every memory strategy, every reader and both benchmarks. \texttt{\{context\}} holds the
memory text built by the strategy (Section~\ref{app:repro_example} shows examples).
\texttt{\{date\_line\}} is empty on LoCoMo; on LongMemEval it gives the question date
(``Today's date is \ldots''), so that relative terms can be resolved. The answer is read from
between the \texttt{<answer>} tags.

\begin{small}
\begin{verbatim}
You are answering a question based on a conversation history.

Conversation history:
{context}

{date_line}Question: {question}
Rules:
1. Answer the question in as few words as possible.
2. For dates, give the date only (e.g., "2022" or "7 May 2023"). Resolve
   relative terms like "yesterday", "last week", "last year" mathematically
   using the timestamp of the fact.
3. For names, the name only.
4. For yes/no questions, answer "Yes" or "No".
5. Do not explain.
6. Do not repeat the question.
7. Put the final answer between <answer> and </answer>.
Answer:
\end{verbatim}
\end{small}

\subsection*{Official LoCoMo prompt (prompt-robustness check only)}

Used only for the substitution test of Section~\ref{sec:validity}. It is reproduced verbatim
from the LoCoMo repository, including its original spelling of ``wriiten''; \texttt{\{s1\}} and
\texttt{\{s2\}} are the two speakers.

\begin{small}
\begin{verbatim}
Below is a conversation between two people: {s1} and {s2}. The conversation
takes place over multiple days and the date of each conversation is wriiten
at the beginning of the conversation.

{context}

Based on the above context, write an answer in the form of a short phrase
for the following question. Answer with exact words from the context
whenever possible.

Question: {question} Short answer:
\end{verbatim}
\end{small}

\subsection*{Graph extraction prompt (\texttt{triple\_json\_v2})}

Sent once per batch of five turns, each rendered as ``speaker: text''. The reply is parsed as
JSON under an extraction budget of 256 tokens (baseline) or 512 tokens (repaired pipeline).

\begin{small}
\begin{verbatim}
Identify factual relationships in this dialogue segment, preserving all
dates, times, locations, and specific details.

Dialogue:
{dialogue}

Respond with a JSON object containing a 'triples' key which is a list of
objects. Each object in 'triples' must have exactly three string keys:
'subject', 'relation', and 'object'.
Do not use pronouns (he, she, it, I) as subjects; resolve them to concrete
names if clear from context.
Capture key attributes of the speakers and any temporal details.
Example output structure:
{"triples": [{"subject": "alice", "relation": "works at", "object":
"north clinic"}, {"subject": "bob", "relation": "graduated in", "object":
"2019"}]}
\end{verbatim}
\end{small}

\subsection*{Judge prompts (LongMemEval \texttt{get\_anscheck\_prompt})}

The judge receives the question, the reference answer and the reader's parsed answer. A reply
beginning with ``yes'' counts as correct and one beginning with ``no'' as incorrect; a reply that
cannot be parsed is set aside for review rather than scored as wrong. LoCoMo answers use the
standard form:

\begin{small}
\begin{verbatim}
I will give you a question, a correct answer, and a response from a model.
Please answer yes if the response contains the correct answer. Otherwise,
answer no. If the response is equivalent to the correct answer or contains
all the intermediate steps to get the correct answer, you should also answer
yes. If the response only contains a subset of the information required by
the answer, answer no.

Question: {question}

Correct Answer: {gold}

Model Response: {pred}

Is the model response correct? Answer yes or no only.
\end{verbatim}
\end{small}

LongMemEval answers use the same standard form except for four question types, as in the
benchmark's own evaluation:
\begin{itemize}
  \item \textbf{Temporal reasoning} adds: ``In addition, do not penalize off-by-one errors for the
        number of days. If the question asks for the number of days/weeks/months, etc., and the
        model makes off-by-one errors (e.g., predicting 19 days when the answer is 18 days), the
        model's response is still correct.''
  \item \textbf{Knowledge update} replaces the equivalence sentences with: ``If the response
        contains some previous information along with an updated answer, the response should be
        considered as correct as long as the updated answer is the required answer.''
  \item \textbf{Single-session preference} supplies a rubric in place of a correct answer: ``Please
        answer yes if the response satisfies the desired response. Otherwise, answer no. The model
        does not need to reflect all the points in the rubric. The response is correct as long as
        it recalls and utilizes the user's personal information correctly.''
  \item \textbf{Abstention} questions ask whether the model ``correctly identifies the question as
        unanswerable''; the model may say that the information is incomplete, or that other
        information is given but the asked information is not.
\end{itemize}

\section{Question sets}
\label{app:repro_questions}

LoCoMo comparisons use 200 questions from the released ten-conversation set, 50 from each of
Categories 1 to 4, taken in turn from each conversation (round-robin). The
LongMemEval subset is a stratified sample of 60 questions, ten from each of the six question
categories, drawn with seed 42 and stored as a fixed file so that every run sees the same
questions in the same order. Question identifiers are assigned per run and depend on the
sample size, so cross-run joins use the stable identifier carried in each record.

Each run records a SHA-256 hash of its loaded question set. The 200 LoCoMo questions hash to
\texttt{8f7200d571bd6391}\ldots, the 60 LongMemEval questions to \texttt{64f5c89412584b98}\ldots{}
and the 20-question phone slice to \texttt{231afbb0f95fac1b}\ldots{} (first 16 hexadecimal digits).

\subsection*{The 20-question phone slice}

The primary phone comparison uses the 20 questions below, drawn by the same round-robin procedure
with five questions per category. Every phone accuracy in Chapter~\ref{ch:on_device_memory} rests
on them. Questions and reference answers are reproduced as they appear in the dataset, including
its spelling errors.

\begin{description}
  \item[Multi-hop (Category 1)]
  \begin{enumerate}
    \item \texttt{conv-26}: What did Caroline research? \emph{Reference:} Adoption agencies
    \item \texttt{conv-30}: What do Jon and Gina both have in common? \emph{Reference:} They lost their jobs and decided to start their own businesses.
    \item \texttt{conv-41}: What martial arts has John done? \emph{Reference:} Kickboxing, Taekwondo
    \item \texttt{conv-42}: What kind of interests do Joanna and Nate share? \emph{Reference:} Watching movies, making desserts
    \item \texttt{conv-43}: what are John's goals with regards to his basketball career? \emph{Reference:} improve shooting percentage, win a championship
  \end{enumerate}
  \item[Temporal reasoning (Category 2)]
  \begin{enumerate}
    \item \texttt{conv-26}: When did Caroline go to the LGBTQ support group? \emph{Reference:} 7 May 2023
    \item \texttt{conv-30}: When Jon has lost his job as a banker? \emph{Reference:} 19 January, 2023
    \item \texttt{conv-41}: Who did Maria have dinner with on May 3, 2023? \emph{Reference:} her mother
    \item \texttt{conv-42}: When did Joanna first watch ``Eternal Sunshine of the Spotless Mind? \emph{Reference:} 2019
    \item \texttt{conv-43}: In which month's game did John achieve a career-high score in points? \emph{Reference:} June 2023
  \end{enumerate}
  \item[Open-domain knowledge (Category 3)]
  \begin{enumerate}
    \item \texttt{conv-26}: What fields would Caroline be likely to pursue in her educaton? \emph{Reference:} Psychology, counseling certification
    \item \texttt{conv-41}: What might John's financial status be? \emph{Reference:} Middle-class or wealthy
    \item \texttt{conv-42}: Is it likely that Nate has friends besides Joanna? \emph{Reference:} Yesteammates on hisvideo game team.
    \item \texttt{conv-43}: Would Tim enjoy reading books by C. S. Lewis or John Greene? \emph{Reference:} C. S.Lewis
    \item \texttt{conv-44}: What is an indoor activity that Andrew would enjoy doing while make his dog happy? \emph{Reference:} cook dog treats
  \end{enumerate}
  \item[Single-hop (Category 4)]
  \begin{enumerate}
    \item \texttt{conv-26}: What did the charity race raise awareness for? \emph{Reference:} mental health
    \item \texttt{conv-30}: How do Jon and Gina both like to destress? \emph{Reference:} by dancing
    \item \texttt{conv-41}: What is John's main focus in local politics? \emph{Reference:} Improving education and infrastructure
    \item \texttt{conv-42}: What is one of Joanna's favorite movies? \emph{Reference:} ``Eternal Sunshineof the Spotless Mind''
    \item \texttt{conv-43}: Which team did John sign with on 21 May, 2023? \emph{Reference:} The Minnesota Wolves
  \end{enumerate}
\end{description}

\section{Memory configuration}
\label{app:repro_memory}

\begin{itemize}
  \item \textbf{Truncated full context.} The most recent turns that fit the reader's context
        budget, which is 4,096 tokens throughout, less a reserve for the prompt template and
        the answer.
  \item \textbf{Vector retrieval.} Turn-level embeddings with mean pooling, cosine similarity,
        and the ten most similar turns by default; Chapter~\ref{ch:memory_evaluation} sweeps
        this from 2 to 50. Retrieved turns are re-sorted into chronological order before they
        enter the prompt. The embedding model expects task prefixes: stored turns are indexed
        as \texttt{search\_document:~} and questions are encoded as \texttt{search\_query:~}.
        Runs that predate this repair are identified where they appear. On the workstation the
        repair raised vector accuracy by 3.5 points for E2B and 6.0 for E4B, and changed
        triple-only graph accuracy by at most half a point.
  \item \textbf{Graph memory.} Triples are extracted from batches of five turns and stored in a
        directed multigraph. Edges start at weight 0.6 and are reinforced by 0.3 on repetition,
        capped at 1.0, with a decay factor of 1.0 and a pruning threshold of 0.15. Retrieval
        returns at most twelve triples. The extraction token budget is 256 in the baseline
        configuration and 512 in the repaired pipeline of Chapter~\ref{ch:graph_mechanisms}.
\end{itemize}

\section{Worked prompt example}
\label{app:repro_example}

The four context formats below come from the recorded Gemma~4~E2B workstation prompts for the
LoCoMo question ``What martial arts has John done?'' (\texttt{conv-41}; reference answer:
kickboxing and taekwondo), which Chapter~\ref{ch:graph_mechanisms} also discusses. Each block
replaces \texttt{\{context\}} in the reader prompt; ``[\ldots]'' marks omitted lines, and long
lines are wrapped for display.

\paragraph{Truncated full context.} 85 turns were sent and the server trimmed the oldest to fit
4,096 tokens. The retained suffix begins in August 2023, after every turn that mentions kickboxing
or taekwondo. Answer: ``John has not mentioned any martial arts.'' (judged incorrect).

\begin{small}
\begin{verbatim}
[5:19 pm on 5 August, 2023] John: I appreciate it! Your words are really
motivating.
[...]
[11:08 am on 16 August, 2023] John: Yeah, Maria, let's keep each other and
everyone else motivated to make a difference! Together, our impact will
surely last.
\end{verbatim}
\end{small}

\paragraph{Vector retrieval, $k=10$.} Ten turns re-sorted into chronological order, 628 prompt
tokens. Answer: ``kickboxing and taekwondo'' (judged correct).

\begin{small}
\begin{verbatim}
[11:01 am on 17 December, 2022] John: Woah, Maria, that sounds cool! I'm
doing kickboxing and it's giving me so much energy.
[6:10 pm on 22 December, 2022] John:  Yep, let's keep making great memories
with our loved ones and cherishing the time we have. I'm off to do some
taekwondo!
[...]
[6:21 pm on 22 July, 2023] John: Yeah, they offer a a bunch, like yoga,
kickboxing, and circuit training. It keeps things interesting!
[...]
\end{verbatim}
\end{small}

\paragraph{Graph triples, version 3.} Twelve selected triples rendered as timestamped statements,
511 prompt tokens. None names either sport. Answer: ``No information available'' (judged
incorrect).

\begin{small}
\begin{verbatim}
[6:03 pm on 6 March, 2023] System: Maria volunteering at Homeless shelter.
[11:08 am on 16 August, 2023] System: John expressed gratitude to Maria.
[...]
[6:21 pm on 22 July, 2023] System: Maria asked about location Martial arts
place or yoga studio.
[6:21 pm on 22 July, 2023] System: Shared photo depicts Two women standing in
a room full of black mats.
[11:01 am on 17 December, 2022] System: John returned from Family road trip
yesterday.
\end{verbatim}
\end{small}

\paragraph{Graph triples with sources, version 4.} The same twelve triples followed by five source
turns, 774 prompt tokens. Answer: ``kickboxing and taekwondo'' (judged correct).

\begin{small}
\begin{verbatim}
[... the twelve triples of version 3 ...]

Source: [11:01 am on 17 December, 2022] John: Woah, Maria, that sounds cool!
I'm doing kickboxing and it's giving me so much energy.
[...]
Source: [6:10 pm on 22 December, 2022] John:  Yep, let's keep making great
memories with our loved ones and cherishing the time we have. I'm off to do
some taekwondo!
[...]
\end{verbatim}
\end{small}

\paragraph{Source-expansion loader correction.}
\label{app:repro_graph_loader}
An early source-restoration run loaded a cached graph without reapplying the source-append
option. The loader was corrected to restore the requested query settings, and that invalid
attempt was excluded from the repair results.

\section{Model artifacts}
\label{app:repro_models}

Every reported cell uses a community-distributed GGUF build at standard
\texttt{Q4\_K\_M} quantization, with uniform per-layer bit widths. The single exception is the
quantization check described in Chapter~\ref{ch:methodology}, which repeated one phone cell on
an adaptive \texttt{UD-Q4\_K\_XL} build and reproduced its verdicts exactly.
Table~\ref{tab:app_artifacts} gives the identity of each file as recovered from the device
cache.

\begin{table}[htbp]
  \centering
  \small
  \begin{tabular}{llrl}
    \toprule
    Role & Repository & Size (GB) & SHA-256 (first 16) \\
    \midrule
    Gemma 4 E2B reader & \texttt{unsloth/gemma-4-E2B-it-GGUF} & 3.11 & not recorded \\
    Gemma 4 E4B reader & \texttt{unsloth/gemma-4-E4B-it-GGUF} & 4.98 & \texttt{85a896a047553e84} \\
    Qwen 3.5 0.8B reader & \texttt{unsloth/Qwen3.5-0.8B-GGUF} & 0.53 & \texttt{bd258782e35f7f45} \\
    Qwen 3.5 2B reader & \texttt{unsloth/Qwen3.5-2B-GGUF} & 1.28 & \texttt{aaf42c8b7c3cab2b} \\
    Qwen 3.5 4B reader & \texttt{unsloth/Qwen3.5-4B-GGUF} & 2.74 & \texttt{00fe7986ff5f6b46} \\
    Vector embedder & \texttt{nomic-ai/nomic-embed-text-v1.5-GGUF} & 0.08 & \texttt{d4e388894e09cf38} \\
    \bottomrule
  \end{tabular}
  \caption{On-device model artifacts. Digests were recomputed from the cached file bytes;
           full values and the repository revisions are kept with the run records. The E2B
           entry is the later of two revisions used during the campaign, discussed in
           Section~\ref{app:repro_limits}.}
  \label{tab:app_artifacts}
\end{table}

The automatic judge is \texttt{llama3.1:70b-instruct-q4\_K\_M}, served on the workstation and
run at temperature 0 with a 256-token answer budget. Judging is performed on the workstation
in every case and is outside the measured phone workload.

\section{Runtime configuration}
\label{app:repro_runtime}

\paragraph{Workstation.} Models are served through Ollama. The reader context window is
4,096 tokens and decoding is greedy at temperature 0. The long-window full-context reruns
(Section~\ref{sec:ch4_long_fc}) raise the window to 131,072 tokens, or 262,144 for Qwen3.5 4B on
LongMemEval-S, and change nothing else.

\paragraph{Smartphone.} Models are served locally through llama.cpp, one inference slot per
server, with an 8,192-token slot context (16,384 for the long-window sittings), five threads
pinned to the performance cluster through a CPU affinity mask, flash attention enabled, the
model's own chat template applied, and the runtime's reasoning budget set to zero. Every
phone sitting therefore runs without reasoning traces. The thinking conditions evaluated in
the workstation experiments were not repeated on the device, so no on-device energy figure
in this thesis describes a thinking configuration. The embedding server runs separately with
mean pooling. Decoding is greedy, which is a deliberate
deviation from the vendor-recommended sampling settings; Section~\ref{sec:validity} reports
the check that compared it against the recommended operating point. Greedy decoding is what
makes the replicate comparison possible: repeated sittings produced byte-identical
predictions, so the variation between them is physical rather than sampling noise.

\section{Software versions and phone server commands}
\label{app:repro_software}

\paragraph{Phone.} The phone ran Android 15 with llama.cpp built from source in Termux. The binary
reports version \texttt{0.1.2-dev} (build 0, commit not recorded) and was built with Clang 21.0.0
for Android aarch64; the source checkout was at commit \texttt{dcad77cc}. Readers were served with
the command below (Gemma~4~E2B shown; the other readers differ only in the model reference), with
\texttt{-c 16384} for the long-window sittings. Network ports are omitted.

\begin{small}
\begin{verbatim}
llama-server -hf unsloth/gemma-4-E2B-it-GGUF:Q4_K_M -c 8192 -np 1
    --cpu-mask f8 -t 5 -fa on --jinja --reasoning-budget 0 --cache-ram 0
\end{verbatim}
\end{small}

Here \texttt{-np 1} gives one inference slot, \texttt{--cpu-mask f8} pins execution to cores 3 to 7
(the prime core and the four performance cores), \texttt{-t 5} sets five threads, \texttt{-fa on}
enables flash attention, \texttt{--jinja} applies the model's own chat template,
\texttt{--reasoning-budget 0} disables reasoning traces and \texttt{--cache-ram 0} disables the
server's prompt cache between requests. The embedder ran as a separate server:

\begin{small}
\begin{verbatim}
llama-server -hf nomic-ai/nomic-embed-text-v1.5-GGUF:Q4_K_M --embeddings
    --pooling mean
\end{verbatim}
\end{small}

The \texttt{-hf} selector resolves the latest repository revision at download time;
Table~\ref{tab:app_artifacts} gives the file digests, and Section~\ref{app:repro_limits} records the
one revision change during the campaign.

\paragraph{Workstation.} Ollama served the readers and the judge; its server version was not
recorded in the run records. The harness's locked Python environment uses Python 3.12.10 with
NumPy 2.4.3, SciPy 1.17.1, NetworkX 3.6.1 and the Ollama client library 0.6.2.

\section{Instrumentation}
\label{app:repro_instrumentation}

Resource sampling runs alongside the harness, requested every 0.5~seconds (realized 1.2 to
3.2~seconds, median per run, because each battery query takes a second or more), and records the battery
charge counter, voltage and current, memory from the kernel's own reporting, and the
device-wide maximum temperature across the available sensors. Sampling starts after the model
servers have loaded, so server start-up is excluded from the measured spans. A wake lock keeps
the processor running with the screen off. Energy is obtained by voltage-weighted Coulomb
counting from the charge counter rather than from a nominal-voltage assumption.

\section{Analysis scripts}
\label{app:repro_scripts}

Every table, figure and statistical test in the thesis is produced from the stored per-run
records by script, not transcribed by hand. Three scripts carry the reported numbers: one
aggregates the accuracy table across runs, one derives the on-device energy metric and is the
sole method of record for it, and one produces the statistical tests. Re-running them on the
stored records regenerates the reported values.

\section{Known limitations of this record}
\label{app:repro_limits}

Four gaps are stated here rather than left for a reader to find.

\begin{enumerate}
  \item \textbf{The E2B artifact changed during the campaign.} The reader was resolved by
        repository and quantization rather than by pinned revision, and the published file was
        updated between sittings. Byte-comparison of the two builds showed identical model
        tensors; the difference is confined to the chat template, in a branch that handles
        tool-call arguments and that these experiments never exercise. Predictions were
        byte-identical across the sittings that span the change, which is what a
        template-only difference predicts. Later runs pin the revision.
  \item \textbf{Run records do not embed the artifact digest.} The link between a result
        directory and the exact model file is reconstructed from the surviving device cache
        and the server logs, not attested inside each run record. The device audit that
        recomputed digests covered five of the six artifacts; the E2B reader had been
        recorded earlier, by revision and file size only, so Table~\ref{tab:app_artifacts}
        leaves its digest blank rather than carrying an unverified value.
  \item \textbf{The primary phone retrieval campaign remains unprefixed.} The 20-question,
        ten-turn retrieval comparison was not rerun with prefixes. The prefixed 100-question
        runs answer the same 20 questions and give the accuracy comparison in
        Section~\ref{sec:retrieval}; no prefixed energy measurement exists for this comparison.
  \item \textbf{The charge counter can stall at the start of a sitting.} Several sittings show
        a frozen counter for the first few minutes with no later catch-up, which under-counts
        energy. The affected sittings and the size of the effect are reported in
        Appendix~\ref{app:on_device_statistics}.
\end{enumerate}

\section{Glossary}
\label{app:repro_glossary}

\begin{description}
  \item[Reader] The model that receives the memory text and writes the answer.
  \item[E2B, E4B] The two Gemma~4 readers, with about 2.3 and 4.5 billion effective parameters.
  \item[Workstation (``the box'')] The Linux host that runs the larger accuracy study through
        Ollama and hosts the judge.
  \item[Sitting] One complete phone run of a workload, unplugged and measured from start to end;
        repeated sittings measure energy variation.
  \item[Truncated full context] The newest dialogue turns that fit a 4,096-token budget, with
        older turns dropped.
  \item[Vector retrieval (Vector RAG)] The $k$ dialogue turns most similar to the question by
        embedding, re-sorted into chronological order; $k=10$ by default.
  \item[SLM-extracted graph] A graph of (subject, relation, object) triples extracted by the reader
        model itself. Figures and tables label its versions GraphRAG v1--v4; it is not Microsoft's
        GraphRAG (Section~\ref{sec:ch6_scope}).
  \item[v1--v4] The four graph versions of Table~\ref{tab:graph_versions}; v3 shows triples only and
        v4 adds the source turns linked to them.
  \item[Source restoration] The v4 step of adding the linked dialogue turns to the selected
        triples.
  \item[Extraction failure] An extraction call whose output the pipeline could not parse; a valid
        empty list is counted separately.
  \item[SPE, SPE-1, SPE-2, $B$] Selective Provenance Expansion, its heuristic and
        representation-specific scoring versions, and its retrieval-context token budget
        ($B=240$ or $370$).
  \item[Oracle] A post hoc upper bound that chooses v3 or v4 per question using the judged labels;
        not an implementable policy.
  \item[Prefixed, unprefixed] Retrieval with or without the embedding model's task prefixes
        (Section~\ref{sec:retrieval}).
  \item[Prefill, decode] Processing the input prompt, and generating the answer token by token.
  \item[KV cache] The stored attention keys and values that let a model reuse earlier tokens during
        generation.
  \item[Correct answers per kilojoule ($G_E$)] Phone-judged correct answers divided by the
        voltage-weighted energy of the same workload; $G_Q$ is the same count per 5,000~mAh full
        charge.
  \item[Parity gate] The predefined rule that decided whether workstation accuracy could stand in
        for phone accuracy.
\end{description}
